\section{Quantum and Quantum-Inspired Graph Models}
\label{sec:quantum}

\subsection{Continuous-time quantum walks}
The continuous-time quantum walk
(CTQW)~\cite{childs2004quantumwalk,venegas2012quantumwalks} is the quantum analogue of
graph diffusion. Replacing the real, dissipative dynamics of the heat equation by the
complex, unitary Schr\"odinger dynamics gives
\begin{equation}
\begin{aligned}
i\,\dot{\boldsymbol{\psi}}=\Lap\boldsymbol{\psi}
\;\;\Longrightarrow\;\;
\boldsymbol{\psi}(t)&=e^{-it\Lap}\boldsymbol{\psi}(0),\\[2pt]
e^{-it\Lap}&=\Eig\,\mathrm{diag}\!\big(e^{-it\lambda_k}\big)\,\Eig^{*}.
\end{aligned}
\label{eq:qwalk}
\end{equation}
The operator $e^{-it\Lap}$ is \emph{unitary} rather than stochastic, so it preserves
the norm of the state and produces interference between paths.

\textit{Interference, intuitively.} In classical diffusion, the contributions of two
walks that reach the same node always add as nonnegative probabilities, so spreading is
a smooth averaging that loses height as it widens. In a quantum walk the contributions
are complex amplitudes that carry a phase, so two walks reaching the same node can add
constructively or cancel. Constructive interference along the graph's geodesics is what
lets probability mass race outward in a coherent front rather than diffuse, and
destructive interference is what produces the oscillatory tails visible in
Fig.~\ref{fig:operators}. The price of this coherence is fragility: any decoherence,
whether physical device noise or, in the classical surrogate, an averaging that washes
out phase, collapses the quantum walk back toward diffusion. This is the same fragility
we see empirically as higher seed-to-seed variance in Section~\ref{sec:casestudy}.

The physically meaningful quantity is the transition probability,
\begin{equation}
\Prop_{\mathrm{qw}}(t)_{ij}=\big|[e^{-it\Lap}]_{ij}\big|^2 ,
\label{eq:qwprob}
\end{equation}
which is again real and row-stochastic, so it plugs into the same layer template as
the heat kernel and stays param-matched with it.

\subsection{The one difference that matters: ballistic spreading}
Interference changes the spreading law. Whereas the heat kernel spreads diffusively,
\eqref{eq:diffusive}, the quantum walk spreads \emph{ballistically}:
\begin{equation}
\Var\big[\,|e^{-it\Lap}|^2\big]\;\sim\;t^2 .
\label{eq:ballistic}
\end{equation}
The probability mass reaches distant nodes quadratically faster in propagation time.

\paragraph*{Worked example, continued} On the same $81$-node path and impulse, the
quantum-walk probability $|e^{-it\Lap}|^2$ has variance exactly $2t^2$: at $t=1$ it is
$2$ (the same as heat), at $t=4$ it is $32$ (four times heat), and at $t=8$ it is $128$
(eight times heat). The two operators agree at $t=1$ and then separate, the quantum
walk pulling ahead as $t^2$ versus $t$ (Fig.~\ref{fig:expA}). The distance frontier
advances linearly in $t$ for the walk but only as $\sqrt{t}$ for diffusion. This is the
entire mechanistic claim, and it is exact, not statistical; the open question is whether
it survives into an end-task gain once weights and a readout are trained on top.

Fig.~\ref{fig:operators} summarizes the whole framework: one Laplacian, three
operators, two spreading laws. This is the single, measurable sense in which the
quantum walk differs from classical diffusion, and Section~\ref{sec:casestudy}
verifies \eqref{eq:diffusive} and \eqref{eq:ballistic} directly.

\begin{figure*}[t]
\centering
\resizebox{\textwidth}{!}{%
\begin{tikzpicture}[node distance=6mm]
% center: the Laplacian
\node[box, minimum width=22mm] (lap) {Graph Laplacian\\$\Lap=\Deg-\Adj$};

% three operators
\node[hot, above right=6mm and 22mm of lap] (gcn) {GCN (local)\\$\hat{\Adj}$, $K$ hops};
\node[hot, right=22mm of lap] (heat) {Heat (global)\\$e^{-t\Lap}$};
\node[hot, below right=6mm and 22mm of lap] (qw) {Quantum walk (global)\\$e^{-it\Lap}$};

\draw[flow] (lap) -- (gcn);
\draw[flow] (lap) -- (heat);
\draw[flow] (lap) -- (qw);

% spreading law annotations
\node[plain, right=20mm of gcn] (gcns) {polynomial,\\$K$-hop horizon};
\node[plain, right=20mm of heat] (heats) {diffusive\\$\Var\sim t$};
\node[plain, right=20mm of qw] (qws) {\textbf{ballistic}\\$\Var\sim t^2$};

\draw[flow] (gcn) -- (gcns);
\draw[flow] (heat) -- (heats);
\draw[flow] (qw) -- (qws);

% mini spreading sketches
\begin{scope}[shift={($(heats.east)+(14mm,0)$)}]
  \draw[steel, thick] plot[domain=-1.2:1.2, samples=40]
       (\x, {0.6*exp(-2*\x*\x)});
  \node[font=\scriptsize, text=ink, above] at (0,0.65) {heat};
\end{scope}
\begin{scope}[shift={($(qws.east)+(14mm,0)$)}]
  \draw[amber!80!black, thick] plot[domain=-1.2:1.2, samples=60]
       (\x, {0.45*(cos(deg(6*\x))*exp(-1.5*\x*\x))+0.15});
  \node[font=\scriptsize, text=ink, above] at (0,0.65) {quantum walk};
\end{scope}
\end{tikzpicture}}
\caption{The unifying view. One Laplacian, three propagation operators, two spreading
laws. The only fundamental difference between classical diffusion and the quantum walk
is diffusive ($\Var\sim t$) versus ballistic ($\Var\sim t^2$) spreading; everything
else in the network is shared, which makes the operators param-matched.}
\label{fig:operators}
\end{figure*}


\subsection{From walks to quantum graph neural networks}
A full quantum graph neural network (QGNN) goes further than a walk. Node features
are encoded into qubits, graph-structured entangling layers (a parameterized quantum
circuit, PQC) play the role of message passing, and measurements form the readout. Four
designs map out the space. \emph{Hamiltonian QGNNs}~\cite{verdon2019qgnn} build the layer
from $e^{-iH_G t}$ for a graph Hamiltonian $H_G$ with trainable couplings, the closest
relative of the quantum walk and the most direct realization of the operator view.
\emph{Quantum-walk neural networks}~\cite{dernbach2019qwnn} make the walk's coin
operator feature-dependent and learnable, so the propagation itself adapts to the data.
\emph{Equivariant quantum graph circuits}~\cite{mernyei2022equivariant} enforce
permutation equivariance in the circuit, the quantum analogue of the symmetry that makes
classical GNNs generalize across node orderings. \emph{Quantum graph
kernels}~\cite{henry2021quantumkernel} skip end-to-end learning and instead use a quantum
evolution to define a similarity between graphs that a classical model then consumes,
which is attractive because it sidesteps barren plateaus. The promise across all four is
twofold: ballistic propagation as above, and an exponentially large Hilbert-space feature
map~\cite{schuld2019featurespace,havlicek2019supervised} that a classical model cannot
represent compactly.
The cost is equally real, and we return to it in Section~\ref{sec:roadmap}: present
quantum hardware is noisy, qubit-limited, and prone to barren plateaus that flatten
the trainable landscape.

\subsection{Encodings and expressivity (a primer)}
A QGNN needs the graph and its features inside a quantum state, and the encoding choice
drives both expressivity and trainability. Three families recur. \emph{Basis encoding}
writes a length-$N$ signal into $\lceil\log_2 N\rceil$ qubits, the convention our
noise study uses, so a walk on $N$ nodes costs only logarithmically many qubits but needs
a unitary that realizes $e^{-it\Lap}$. \emph{Angle encoding} maps each feature to a
rotation angle, one qubit per feature, and is the workhorse of variational quantum
circuits~\cite{mitarai2018qcl,farhi2018classification}. \emph{Amplitude encoding} packs a
normalized feature vector into the amplitudes of a state, exponentially compact but
expensive to prepare. Table~\ref{tab:encodings} contrasts the three.

\begin{table}[t]
\centering
\caption{Quantum feature encodings for a length-$N$ signal.}
\label{tab:encodings}
\renewcommand{\arraystretch}{1.25}
\footnotesize
\begin{tabular}{@{}p{0.20\columnwidth}p{0.16\columnwidth}p{0.18\columnwidth}p{0.20\columnwidth}@{}}
\toprule
Encoding & Qubits & Prep cost & Typical use \\
\midrule
Basis & $\lceil\log_2 N\rceil$ & low (a basis state) & walks, search (our Exp~D) \\
Angle & $O(N)$ & low (one rotation each) & variational circuits \\
Amplitude & $\lceil\log_2 N\rceil$ & high (state prep) & compact feature maps \\
\bottomrule
\end{tabular}
\end{table} The promise of these encodings is a feature map into a Hilbert
space whose dimension grows exponentially with the qubit
count~\cite{schuld2019featurespace,havlicek2019supervised,schuld2018book}, where a linear
decision boundary can separate data that is hard to separate classically, and in
restricted settings this yields provable data or sample
advantages~\cite{huang2021power,liu2021rigorous,abbas2021power}. The catch is that the
same expressivity makes the training landscape flat: random deep circuits exhibit barren
plateaus where gradients vanish exponentially in the number of
qubits~\cite{mcclean2018barren}, so expressivity and trainability trade off. For a
practitioner the takeaway is that a quantum graph model is not automatically better; its
advantage, if any, depends on a careful encoding and a problem whose structure the
encoding actually exploits, which is precisely what a controlled, honest evaluation is
meant to detect.

\subsection{Quantum-inspired: the runnable middle ground}
Between a classical GCN and a fault-tolerant QGNN sits the option we recommend for
near-term practice. The quantum-walk kernel \eqref{eq:qwprob} is just a matrix
function of the Laplacian; it can be computed on a classical processor with one
eigendecomposition and used as a propagation layer. This \emph{quantum-inspired}
GNN keeps the ballistic spreading mechanism while remaining trainable and
reproducible today. It is the operator we build in the next section and evaluate in
the case study. Whether its ballistic mechanism translates into an end-task advantage
is precisely the empirical question we refuse to answer by assertion.
